\documentclass[11pt]{article}
\usepackage[a4paper,margin=1in]{geometry}
\usepackage{fontspec}
\usepackage[boldfont=false]{xeCJK}
\setCJKmainfont{FandolSong-Regular.otf}
\usepackage{amsmath}
\usepackage{amssymb}
\usepackage{array}
\usepackage{booktabs}
\usepackage{graphicx}
\usepackage{adjustbox}
\usepackage{enumitem}
\setlist[itemize]{topsep=2pt,partopsep=0pt,parsep=0pt,itemsep=2pt,leftmargin=1.4em}
\setlist[enumerate]{topsep=2pt,partopsep=0pt,parsep=0pt,itemsep=2pt,leftmargin=1.6em}
\usepackage{parskip}
\usepackage{fvextra}
\fvset{breaklines=true,breakanywhere=true,fontsize=\small}
\setlength{\parindent}{0pt}

\begin{document}

\begin{center}
  {\LARGE\bfseries Government macroeconomic intervention}\\[0.35em]
  {\large A Level Economics}
\end{center}
\vspace{0.6em}

Governments try to steer the whole economy. This handout covers what they aim for and the three sets of tools they use.

\section*{Macroeconomic policy objectives}

A government usually has five main aims:

\begin{itemize}
\item \textbf{price stability} 物价稳定 — low and steady \textbf{inflation} 通货膨胀 (often a target of about 2\% a year).

\item \textbf{full employment} 充分就业 — jobs for almost everyone who wants to work.

\item \textbf{economic growth} 经济增长 — a rising real output over time.

\item \textbf{balance of payments} 国际收支 stability — not importing far more than the country exports for years on end.

\item \textbf{redistribution of income} 收入再分配 — narrowing the gap between rich and poor.

\end{itemize}

\subsection*{Conflicts between objectives}

The aims can \textbf{conflict} 冲突 — reaching one can hurt another:

\begin{itemize}
\item \textbf{growth versus inflation}: fast growth raises \textbf{aggregate demand} 总需求, which can pull prices up.

\item \textbf{growth versus the balance of payments}: as incomes rise, people buy more imports, worsening the trade balance.

\item \textbf{unemployment versus inflation}: cutting unemployment by raising demand can cause more inflation.

\item \textbf{growth versus the environment}: more output often means more pollution.

\end{itemize}

The government has three sets of tools. Fiscal and monetary policy are \textbf{demand-side policies} 需求侧政策 (they work on AD). Supply-side policy works on \textbf{aggregate supply} 总供给.

\section*{Fiscal policy}

\textbf{fiscal policy} 财政政策 is the use of government \textbf{spending} and \textbf{taxation} to influence the economy.

The \textbf{budget} 预算 compares the two:

\begin{itemize}
\item a \textbf{budget deficit} 预算赤字 is when spending is greater than tax revenue. The government must borrow, adding to the \textbf{national debt} 国债.

\item a \textbf{budget surplus} 预算盈余 is when tax revenue is greater than spending.

\end{itemize}

\subsection*{Types of tax}

\begin{itemize}
\item a \textbf{direct tax} 直接税 is paid straight to the government on income or profit (income tax, corporation tax).

\item an \textbf{indirect tax} 间接税 is a tax on spending (added to the price of goods).

\end{itemize}

Taxes are also grouped by how the rate changes with income:

\begin{center}
\adjustbox{max width=\linewidth}{%
\begin{tabular}{ll}
\toprule
\textbf{Type} & \textbf{How it works} \\
\midrule
\textbf{progressive tax} 累进税 & the rate \textbf{rises} as income rises (the rich pay a larger share) \\
\textbf{proportional tax} 比例税 & the rate \textbf{stays the same} at every income \\
\textbf{regressive tax} 累退税 & the rate \textbf{falls} as income rises (the poor pay a larger share) \\
\bottomrule
\end{tabular}}
\end{center}

\subsection*{Using fiscal policy}

\begin{itemize}
\item \textbf{expansionary} 扩张性 fiscal policy raises government spending or cuts taxes. This raises AD, to fight a recession and unemployment.

\item \textbf{contractionary} 紧缩性 fiscal policy cuts spending or raises taxes. This lowers AD, to fight inflation.

\end{itemize}

\textbf{Effectiveness.} Fiscal policy can act fast on AD, but it has problems: time lags (a budget is set once a year), large deficits and debt, and \textbf{crowding out} 挤出效应 — when heavy government borrowing pushes up interest rates and reduces private investment.

\section*{Monetary policy}

\textbf{monetary policy} 货币政策 is run by the \textbf{central bank} 中央银行. Its main tool is the \textbf{interest rate} 利率, but it can also change the \textbf{money supply} 货币供给 or influence the \textbf{exchange rate} 汇率.

\subsection*{The transmission mechanism}

The \textbf{transmission mechanism} 传导机制 is how a change in the interest rate reaches the wider economy. If the central bank \textbf{raises} the interest rate:

\begin{itemize}
\item borrowing becomes dearer, so households and firms spend and invest less.

\item saving is rewarded, so people spend less.

\item the currency tends to rise, so exports fall.

\end{itemize}

All of these lower AD, which slows inflation. A \textbf{cut} in the interest rate does the opposite, raising AD.

\textbf{Effectiveness.} Monetary policy is flexible (the rate can change at any time), but it works with a time lag of a year or more. In a deep slump, even very low rates may not make worried firms and households borrow.

\section*{Supply-side policy}

\textbf{supply-side policy} 供给侧政策 tries to raise the economy's \textbf{productive capacity} 生产能力 — to shift long-run aggregate supply to the right. It aims to make markets work better and raise \textbf{productivity} 生产率.

There are two kinds.

\textbf{Market-based} policies let markets work more freely:

\begin{itemize}
\item cutting income tax and benefits to strengthen the \textbf{incentive} 激励 to work.

\item \textbf{privatisation} 私有化 — selling state firms to private owners, who run them for profit.

\item \textbf{deregulation} 放松管制 — removing rules that block competition.

\end{itemize}

\textbf{Interventionist} policies use the government directly:

\begin{itemize}
\item spending on education and \textbf{training} 培训 to raise workers' skills.

\item spending on \textbf{infrastructure} 基础设施 — roads, ports and broadband.

\item support for research and new technology.

\end{itemize}

\textbf{Effectiveness.} Supply-side policy is powerful because it can raise growth \textbf{and} lower inflation at the same time, with no trade-off. But it is slow (education takes years to pay off), costly, and some measures, like cutting benefits, can widen inequality.


\end{document}
